% !TeX root = ../教材题库_全部条目_转写工作稿.tex
% Source: XJ-B1-C02-EX-P052-Y0592
% 第2章 / 2.3 一元二次不等式 / 教材 p.52 / PDF 第59页
% OCR未校对：公式与文字均可能有误；未完成数学TeX转写；禁止直接发给学生
% 原题及相邻原解截图（核对用）：../题目原图/XJ-B1-C02-EX-P052-Y0592-01.webp
% 本文件只在转写完成后作为题干使用；OCR 条目尚不是可靠题干。
\begin{bankitem}{2.3 例题 例题1}{第2章 / 2.3 一元二次不等式 / 教材 p.52 / PDF 第59页}{XJ-B1-C02-EX-P052-Y0592}{OCR未校对：公式与文字均可能有误；未完成数学TeX转写；禁止直接发给学生}
\begin{ocrdraft}
[教材p.52 / PDF第59页]
\par
00解不等式2亠一3冫0．
\par
解由于方程22亠一3一0有两个不相等的实数根
\par
3
\par
则函数3，一2—x—3的图象与x轴有两个交点（一1，
\par
，如图2．3
\par
5所示．
\par
由图象得不等式22亠一3冫0的解集为
\par
3
\par
一1或冫
\par
2
\par
y=2x2—x—3
\par
3
\par
y=4x—4x+l
\par
2
\par
0
\par
图2．3一5
\par
图2．3一6
\end{ocrdraft}
\end{bankitem}
